%%%%%%%%%%%%%%%%%%%%%%%%%%%%%%%%%%%%%%%%%
% Medium Length Professional CV
% LaTeX Template
% Version 3.0 (December 17, 2022)
%
% This template originates from:
% https://www.LaTeXTemplates.com
%
% Author:
% Vel (vel@latextemplates.com)
%
% Original author:
% Trey Hunner (http://www.treyhunner.com/)
%
% License:
% CC BY-NC-SA 4.0 (https://creativecommons.org/licenses/by-nc-sa/4.0/)
%
%%%%%%%%%%%%%%%%%%%%%%%%%%%%%%%%%%%%%%%%%

%----------------------------------------------------------------------------------------
%	PACKAGES AND OTHER DOCUMENT CONFIGURATIONS
%----------------------------------------------------------------------------------------

\documentclass[
	%a4paper, % Uncomment for A4 paper size (default is US letter)
	11pt, % Default font size, can use 10pt, 11pt or 12pt
]{resume} % Use the resume class

\usepackage{ebgaramond} % Use the EB Garamond font
\usepackage[dvipdfm,linkcolor=blue,anchorcolor=blue,urlcolor=blue,colorlinks=true]{hyperref}
\urlstyle{same}
\usepackage{pxjahyper}

\name{研究業績} % Your name to appear at the top

\address{\parbox{\textwidth}{\raggedleft \Large \bf{日下部 完}}}
\address{\parbox{\textwidth}{\raggedleft 北海道大学大学院情報科学院 \\情報科学専攻 情報理工学コース \\ヒューマンコンピュータインタラクション研究室}}
\address{\parbox{\textwidth}{\raggedleft \url{kusakabe.kan.v5@elms.hokudai.ac.jp}}}
\address{\parbox{\textwidth}{\raggedleft \url{https://kankusakabe.github.io}}}


\begin{document}



\begin{rSection}{査読付き国際会議}{}{}{}
		\item Yuki Abe*, \underline{\textbf{Kan Kusakabe*}}, Myungguen Choi*, Daisuke Sakamoto, Tetsuo Ono. *Joint first authors. Understanding Usability of VR Pointing Methods with a Handheld-style HMD for Onsite Exhibitions. ACM CHI Conference on Human Factors in Computing Systems (CHI 2025), ACM, April 2025. (採択率: 27 \%; \textbf{Honorable Mentions: top 5 \% of accepted papers}). [\href{https://doi.org/10.1145/3706598.3713874}{doi}]
\end{rSection}

\begin{rSection}{査読付き論文誌}{}{}{}
	\item \underline{\textbf{日下部 完}}, 坂本 大介, 小野 哲雄. RingSense：スマートフォン背面のジェスチャ入力を実現するスマホリング型インタフェースの設計と実装 情報処理学会論文誌，情報処理学会, Vol. 66 (2). [\href{https://doi.org/10.20729/0002000028}{doi}]
	\item \underline{\textbf{日下部 完}}, 坂本 大介, 小野 哲雄. BodSockets：スマホソケット型インタフェースを用いたスマートフォンの背面領域を活用した操作手法の提案 情報処理学会論文誌，情報処理学会, Vol. 66 (12). (採録予定)
\end{rSection}

\begin{rSection}{査読付き国内会議}{}{}{}
		\item \underline{\textbf{日下部 完}}, 坂本 大介, 小野 哲雄. スマートフォン背面のジェスチャ入力を実現するスマホリング型デバイスの設計と実装. 第30回インタラクティブシステムとソフトウェアに関するワークショップ（WISS 2022），日本ソフトウェア科学会, pp.8-14，December 2022. (\textbf{Long-press, acceptance rate: 12.5\%}). [\href{https://www.wiss.org/WISS2022Proceedings/data/02.pdf}{link}]
\end{rSection}

\begin{rSection}{その他}{}{}{}
	\item \underline{\textbf{日下部 完}}*, 阿部優樹*, 坂本 大介, 小野 哲雄. Game-2-X: 種類が異なるゲームプレイ間を繋ぐシステムの提案. （* 共同第一著者）第31回インタラクティブシステムとソフトウェアに関するワークショップ（WISS 2023），日本ソフトウェア科学会, デモンストレーション発表. December 2023. [\href{https://www.wiss.org/WISS2023Proceedings/data/1-A03.pdf}{link}]
	\item \underline{\textbf{日下部 完}}, 崔 明根, 坂本 大介, 小野 哲雄. 無段階調整インタフェースのためのハンドジェスチャによる操作手法の探索的研究. 情報処理学会 研究報告ヒューマンコンピュータインタラクション（HCI）, 2022-HCI-197(25), 1-8, March 2022. [\href{http://id.nii.ac.jp/1001/00217337/}{link}]
\end{rSection}

\begin{rSection}{招待講演}{}{}{}
	\item 市民公開講座 / 横浜市MICE次世代育成事業 CHI 2025シンポジウム 触れる未来、創るHCIの世界. April 26 2025. [\href{https://sigchi.jp/symposium/chi2025_symposium/}{url}]
	\item FIT2025 第24回情報科学技術フォーラム トップコンファレンスセッション. Sep 5 2025. [\href{https://www.gakkai-web.net/fit/program_web/data/html/event/TCS6-1.html}{url}]
\end{rSection}

\begin{rSection}{競合的資金}
1. Hokkaido University EXEX Doctoral Fellowship. 2024--2027. \$3,000 annual and stipend.

2. HOKUDAI TECH GARAGE Project. 2 months (March-April 2022), \$ 2,000 per project.
\end{rSection}







\end{document}
